\section{A certificate for one point of the diagram}\label{cp:sec:cert}

This section writes out the finite part of the statement $(0.2,0.6)\in\mathcal G$ as a certified computation with explicit constants, and runs it in ball arithmetic below height $160$. The statement is that $\xi_{0.2}$ has no zero with real part above $0.6$. Given a barrier $X$ (in the variable of $H_{0.2}$, so that it sits at height $X/2$ for $\xi_{0.2}$; see Theorem~\ref{cp:thm:cert}) beyond which this is known \cite{Polymath}, it is the vanishing of the winding number of $\xi_{0.2}$ around the rectangle $Q=[0.6,0.9]\times[0,X/2]$. Throughout, $t=0.2$.

\subsection{The lemmas}\label{cp:sec:lemmas}

\begin{lemma}[Bottom edge]\label{cp:lem:bottom}
$\xi_t(\sigma)>0$ for every real $\sigma$ and every $t\ge0$.
\end{lemma}

\begin{proof}
Lemma~\ref{cp:lem:structure}(ii).
\end{proof}

\begin{lemma}[Right edge]\label{cp:lem:right}
For $0\le t\le\frac12$ every zero of $\xi_t$ satisfies $|\beta-\frac12|\le\frac12\sqrt{1-2t}$. In particular for $t=0.2$ every zero has $\beta\le0.8873$, so $\xi_{0.2}$ has no zero on $\re s=0.9$.
\end{lemma}

\begin{proof}
De Bruijn's theorem on strips, as in the proof of Theorem~\ref{cp:thm:phase}, with $\Delta=1$ from the critical strip and $\beta-\frac12=\frac12\im z$ up to sign.
\end{proof}

\begin{lemma}[Tilted bound]\label{cp:lem:tilt}
Let $0<\eta<\pi/8$, $t\ge0$ and $s=\frac12+x+i\theta$ with $\theta\ge0$. Then
\[
|\xi_t(s)|\le K_\eta(t,x)\,e^{-2\eta\theta},\qquad K_\eta(t,x)=4\int_{-\infty}^{\infty}\big|e^{t(u+i\eta)^2}\big|\,|\Phi_{\rm dBN}(u+i\eta)|\,e^{2xu}\,du,
\]
and $K_\eta(t,x)$ is finite and even in $x$. Write $\bar K_\eta(t,R)=\max_{|x|\le R}K_\eta(t,x)$.
\end{lemma}

\begin{proof}
Since $\Phi_{\rm dBN}$ is even, $\xi_t(s)=4\int_\R e^{tu^2}\Phi_{\rm dBN}(u)e^{2(x+i\theta)u}\,du$. The kernel extends analytically to $|\im u|<\pi/8$ with the same doubly exponential decay, which is Proposition~\ref{prop:tilt} in the variable $v=2u$, so the contour may be moved to $\im u=\eta$, and $|e^{2i\theta(u+i\eta)}|=e^{-2\eta\theta}$.
\end{proof}

\begin{lemma}[Derivative bound]\label{cp:lem:deriv}
With the notation of Lemma~\ref{cp:lem:tilt}, for $r>0$, $|\xi_t'(s)|\le r^{-1}\bar K_\eta(t,|x|+r)\,e^{-2\eta(\theta-r)}$, and the factor $e^{-2\eta(\theta-r)}$ may be replaced by $1$ when $\theta\le r$.
\end{lemma}

\begin{proof}
Cauchy's estimate on the disc of radius $r$ about $s$. A point $z$ of the disc has $|\re z-\frac12|\le|x|+r$ and $\im z\ge\theta-r$. If $\im z\ge0$, Lemma~\ref{cp:lem:tilt} bounds $|\xi_t(z)|$ by $\bar K_\eta(t,|x|+r)e^{-2\eta\im z}$, which is at most $\bar K_\eta(t,|x|+r)e^{-2\eta(\theta-r)}$ and at most $\bar K_\eta(t,|x|+r)$. If $\im z<0$, which occurs only when $\theta<r$, then $\xi_t(\bar z)=\overline{\xi_t(z)}$ gives the bound $\bar K_\eta(t,|x|+r)$ for $|\xi_t(z)|$. So on the disc $|\xi_t|\le\bar K_\eta(t,|x|+r)e^{-2\eta(\theta-r)}$, and $|\xi_t|\le\bar K_\eta(t,|x|+r)$ when $\theta\le r$.
\end{proof}

\begin{lemma}[Grid]\label{cp:lem:grid}
Let $s_0,s_1$ be points of a segment on which $\xi_t$ has no zero, and suppose $|s_1-s_0|\cdot\sup_{[s_0,s_1]}|\xi_t'|<|\xi_t(s_0)|$. Then the change of $\arg\xi_t$ along the segment equals the principal value $\operatorname{Arg}(\xi_t(s_1)/\xi_t(s_0))\in(-\pi,\pi)$.
\end{lemma}

\begin{proof}
The image of the segment lies in the open disc of radius $|\xi_t(s_0)|$ about $\xi_t(s_0)$, which does not contain $0$, and on that disc a continuous branch of the argument varies by less than $\pi$ from its value at the centre.
\end{proof}

\begin{theorem}[The certificate]\label{cp:thm:cert}
Let $X$ be a barrier for $(0.2,0.6)$ in the variable of $H_{0.2}$, as Polymath15 provide one \cite{Polymath} (no script of this book checks it), that is, a number such that $H_{0.2}$ has no zero $z$ with $\re z\ge X$ and $\im z>0.2$, so that $\xi_{0.2}$ has no zero with $\re s>0.6$ and $\im s\ge X/2$. Let $s_0,\dots,s_M$ be points of the left, top and right edges of $Q$, in order along $\partial Q$ from one bottom corner, $s_0$, to the other, $s_M$, and including the two top corners, and let $B_j$ be ball enclosures of $\xi_{0.2}(s_j)$ with radii $\varepsilon_j$. Suppose that for every $j$ there are $r>0$ and $\eta\in(0,\frac\pi8)$ such that
\[\begin{gathered}
|s_{j+1}-s_j|\,\frac{\bar K_\eta(0.2,\,0.4+r)}{r}\,e^{-2\eta(\theta_j-r)}+\varepsilon_j+\varepsilon_{j+1}<\min|B_j|,\\ \theta_j=\min(\im s_j,\im s_{j+1}),
\end{gathered}\]
where the factor $e^{-2\eta(\theta_j-r)}$ may be replaced by $1$ when $\theta_j\le r$, and that $\sum_j\operatorname{Arg}(B_{j+1}/B_j)$, evaluated in ball arithmetic, lies within $(-\pi,\pi)$. Then $\xi_{0.2}$ has no zero in $Q$, so $(0.2,0.6)\in\mathcal G$ and $\Lambda\le0.22$.
\end{theorem}

\begin{proof}
By Lemma~\ref{cp:lem:bottom}, $\xi_{0.2}$ is real and positive on the bottom edge, so $\arg\xi_{0.2}$ does not change along it, and by Lemma~\ref{cp:lem:right} the right edge is zero free. The hypothesis gives $|\xi_{0.2}(s_j)|\ge\min|B_j|$ at the grid points, and Lemma~\ref{cp:lem:deriv} with $|x|\le0.4$ and $\im z\ge\theta_j$ on the segment from $s_j$ to $s_{j+1}$ bounds the derivative between them, so Lemma~\ref{cp:lem:grid} applies to every step, and the winding number is the integer nearest to $\frac1{2\pi}\sum_j\operatorname{Arg}(B_{j+1}/B_j)$, which is $0$. So $\xi_{0.2}$ has no zero in $Q$; it has none with $\re s>0.9$ by Lemma~\ref{cp:lem:right}, none with $\re s>0.6$ and $\im s\ge X/2$ by the barrier, and its zeros in $\im s<0$ are the conjugates of those in $\im s>0$. Hence it has no zero with $\re s>0.6$, which is $(0.2,0.6)\in\mathcal G$, and the last clause is Theorem~\ref{cp:thm:phase}.
\end{proof}

The number of grid points is of order $(X/2)\sup_jK_\eta e^{2\eta r}/(r\,|\xi_t(s_j)|e^{2\eta\theta_j})$. The decay of $\xi_t$, of order $e^{-\pi\theta/4}$ away from its zeros, cancels against the decay $e^{-2\eta\theta}$ of the derivative bound only at $\eta=\pi/8$, which Lemma~\ref{cp:lem:tilt} excludes, and $K_\eta$ grows without bound as $\eta\uparrow\pi/8$ (Table~\ref{cp:tab:K}). With a tilt $\eta<\pi/8$ the count grows like $e^{(\pi/4-2\eta)\theta}$. We have no proof that the count is $O(X\log X)$, or polynomial in $X$, when the tilt is chosen at each height. In the runs below the left edge needs on average about $45$ steps per unit height up to height $60$, $73$ up to $120$ and $94$ up to $160$ (Theorem~\ref{cp:thm:rig}).

\begin{proposition}[The certificate in Thorin form]\label{cp:prop:thorincert}
Let $t\ge0$ and $\frac12<\alpha<\sigma_1$, and suppose $v_t(\sigma,\theta)>0$ for every $(\sigma,\theta)\in[\alpha,\sigma_1]\times[0,T]$ at which $\xi_t(\sigma+i\theta)\ne0$. Then $\xi_t$ has no zero $\beta+i\gamma$ with $\alpha<\beta\le\sigma_1$ and $0<\gamma<T$.
\end{proposition}

\begin{proof}
Near a zero $\rho$ of multiplicity $m$, $\xi_t'/\xi_t(s)=m/(s-\rho)+O(1)$, so at $s=\rho-\delta$ with small $\delta>0$ the real part is $-m/\delta+O(1)<0$. If $\alpha<\beta\le\sigma_1$, these points lie in the rectangle for small $\delta$. This is the argument of Remark~\ref{rem:ch2:oneline}, restricted to a rectangle of finite height.
\end{proof}

Positivity on the single line $\re s=\alpha$ is not enough (Remark~\ref{rem:ch2:oneline}), which is why the proposition asks for the whole rectangle. It is the pointwise counterpart of Theorem~\ref{cp:thm:cert}: the same evaluations of $\xi_t$ and $\xi_t'$ certify it, with a two dimensional grid and a bound on the derivative of $v_t$ in place of the argument increments, and the margin is the one measured in Proposition~\ref{cp:prop:see}.

\subsection{Constants and the march}\label{cp:sec:march}

Table~\ref{cp:tab:K} gives the constants of Lemma~\ref{cp:lem:tilt} at $t=0.2$ and the values of $\xi_{0.2}$ on the real axis that fix the scale.

\begin{table}[ht]
\centering\small
\begin{tabular}{@{}lccc@{}}
\toprule
$\eta$ & decay $e^{-2\eta\theta}$ & $K_\eta(0.2,0.4)$ & $K_\eta(0.2,0.9)$\\
\midrule
$0.30$ & $0.60$ & $23.67$ & $25.31$\\
$0.35$ & $0.70$ & $199.2$ & $231.1$\\
$0.38$ & $0.76$ & $4847$ & $6776$\\
\bottomrule
\end{tabular}
\qquad
\begin{tabular}{@{}lc@{}}
\toprule
$\sigma$ & $\xi_{0.2}(\sigma)$\\
\midrule
$0.5$ & $0.49827$\\
$0.6$ & $0.49839$\\
$0.9$ & $0.50013$\\
$1.0$ & $0.50117$\\
$1.5$ & $0.50996$\\
\bottomrule
\end{tabular}
\caption{Constants for the certificate at $t=0.2$, in floating point. Script \texttt{phase\_cert.py}.}
\label{cp:tab:K}
\end{table}

A floating point march of Lemma~\ref{cp:lem:grid} around $Q_{60}=[0.6,0.9]\times[0,60]$, with the P\'olya route evaluator (the trapezoid rule of Section~\ref{sec:ch4:trapezoid} with the weight $e^{tu^2}$) at $40$ digits (working precision $72$ digits) and an adaptive step keeping $|\xi_t(s_{j+1})-\xi_t(s_j)|<\frac12|\xi_t(s_j)|$, returned the winding $-2.9\times10^{-72}$ in $910$ evaluations. The same evaluator located the zeros of $\xi_{0.2}$ on the critical line below $60$,
\[\begin{gathered}
14.077,\ 20.947,\ 24.952,\ 30.330,\ 32.888,\ 37.503,\\ 40.833,\ 43.282,\ 47.887,\ 49.730,\ 52.905,\ 56.358,\ 59.237,
\end{gathered}\]
thirteen zeros, as for $\xi$ below $60$, each moved slightly by the flow (five of them are in Table~\ref{tab:ch21:flowzeros}). For the control function that moves the pair at $30.330$ and $32.888$ to the quadruple $\frac12\pm\frac14+31.609i$ and its conjugates, as in Theorem~\ref{thm:ch13:perturb} (there for $\xi$ and ordinates $1000<A<B\le2000$), the computed winding was $1$ to within $10^{-35}$ in $905$ evaluations: exactly one of the four new zeros, $0.75+31.609i$, lies in $Q_{60}$.

The finite part of Theorem~\ref{cp:thm:cert} was then run in ball arithmetic, adding every truncation error to the radius of the enclosure. The evaluator is the same trapezoid rule, with $w=s-\frac12=x+i\theta$ and $u_j=jh$: $\xi_t(s)=8h\sum_{j=0}^{J}\omega_je^{tu_j^2}\Phi_{\rm dBN}(u_j)\cosh(2wu_j)+E$, where $Jh\ge3$, $\omega_0=\frac12$, $\omega_j=1$ for $j\ge1$, and $|E|\le E_{\rm cut}+E_{\rm trap}$. For $u\ge3$ the kernel satisfies $|\Phi_{\rm dBN}(u)|\le60\,e^{9u}e^{-\pi e^{4u}}$, so that for $|x|\le1$ the summand is at most $8h\,G(u_j)$ with $G(u)=60\exp(tu^2+11u-\pi e^{4u})$, which decreases with $-G'/G\ge\lambda=4\pi e^{4Jh}-2tJh-11$ for $u\ge Jh$, and the nodes beyond $Jh$ contribute at most $E_{\rm cut}=8G(Jh)/\lambda<10^{-229000}$. The error $E_{\rm trap}$ of the rule is bounded by Poisson summation on the strip $|\im u|<a$, $a=0.30<\pi/8$, in which the integrand of Lemma~\ref{cp:lem:tilt} is analytic: $E_{\rm trap}\le2K_{0.30}(0.2,x)\cosh(2a\theta)/(e^{2\pi a/h}-1)$, with $K_{0.30}(0.2,x)\le23.72$ for $|x|\le0.4$, bounded in the way described below, and $h$ is taken so that this is below $5\times10^{-4}e^{-\pi T/4}$ at the top of the rectangle of height $T$: $h=0.0199$, $0.0106$ and $0.008$ at $T=60$, $120$ and $160$, at the working precisions $219$, $288$ and $334$ bits. The series for the kernel is summed until its tail is below $2^{-{\rm prec}-30}$, the tail after $N$ terms being enclosed by
\[\begin{aligned}
&\sum_{n>N}\big|(2\pi^2n^4e^{9u}-3\pi n^2e^{5u})e^{-\pi n^2e^{4u}}\big|\\&\qquad\le|e^{9u}|\Big(2\pi^2+\frac{3\pi}\kappa\Big)\Big(\frac4{\pi\kappa}\Big)^2e^{-2}\,\frac{e^{-\pi(N+1)^2\kappa/2}}{1-e^{-\pi(N+1)\kappa}},\qquad \kappa=\re(e^{4u}),
\end{aligned}\]
using $|e^{5u}|\le|e^{9u}|/\kappa$, $n^4e^{-\pi n^2\kappa/2}\le(4/\pi\kappa)^2e^{-2}$ and $n^2\ge(N+1)^2+2(N+1)(n-N-1)$ for the geometric sum, for every $N$ and every $\kappa>0$. At a real node $\kappa=e^{4u}$; at a complex $u$ $\kappa$ is a lower bound of $\re(e^{4u})$ over the input ball. The constants $\bar K_\eta(0.2,0.4+r)$ were bounded above by interval subdivision of $[0,3]$ into $3\cdot2^{14}$ cells, plus the tail $u\ge3$, below $e^{-25000}$ since $\re e^{4(u+i\eta)}\ge e^{12}\cos(4\eta)>8000$ there, for the three tilts $\eta=0.30$, $0.35$ and $0.38$ and the three radii $r=\frac14$, $\frac38$ and $\frac12$: $24.37$, $24.82$ and $25.36$, then $212.33$, $221.11$ and $231.71$, then $5633.6$, $6170.8$ and $6841.2$, each rounded up. The ordinates of the grid points are exact dyadic numbers, and the abscissae $0.6$ and $0.9$ are balls of radius $2^{-{\rm prec}}$ enclosing them, so that the radius of $\xi_t(s_j)$ is that of the evaluator. At each step the tilt $\eta$ and the radius $r$ giving the smallest derivative bound are used, the step is $0.8\,|\xi_t(s_j)|_{\rm low}$ divided by that bound, capped at $0.5$, and the inequality of Theorem~\ref{cp:thm:cert}, with the exact length of the step and its exact lowest height, is then verified in ball arithmetic, the step being halved when it fails.

\begin{theorem}\label{cp:thm:rig}
$\xi_{0.2}$ has no zero $s$ with $0.6\le\re s\le0.9$ and $0\le\im s\le160$.
\end{theorem}

\begin{proof}
For $T=60$, $120$ and $160$ the script walks the boundary of $Q_T=[0.6,0.9]\times[0,T]$ counterclockwise, up the right edge, along the top and down the left edge, with the steps of Theorem~\ref{cp:thm:cert}. The bottom edge is not walked, since $\xi_{0.2}$ is real and positive on it (Lemma~\ref{cp:lem:bottom}). The run took $4483$ steps for $T=60$ ($1737$ up the right edge, $31$ along the top and $2715$ down the left edge), $13{,}679$ for $T=120$ ($4955$, $11$ and $8713$) and $23{,}219$ for $T=160$ ($8142$, $20$ and $15{,}057$), and returned the winding number enclosures $[\pm3.94\times10^{-5}]$, $[\pm1.24\times10^{-6}]$ and $[\pm3.47\times10^{-7}]$. In every run the inequality of Theorem~\ref{cp:thm:cert} was verified in ball arithmetic at every step and no step failed, and the sum of the arguments lies within $\pi$ of $0$, so by Lemmas~\ref{cp:lem:bottom} to~\ref{cp:lem:grid} the winding number is $0$ and $\xi_{0.2}$ has no zero in $Q_T$, which for $T=160$ is the theorem. Script \texttt{phase\_rig.py}, data \texttt{phase\_rig.json}.
\end{proof}

With Lemma~\ref{cp:lem:right} this says that every zero of $\xi_{0.2}$ below height $160$ has $|\beta-\frac12|<0.1$; as a control, the same script run around $[0.4,0.6]\times[0,60]$, which contains the thirteen zeros on the line, returns the winding number enclosure $[13.0000\pm5.04\times10^{-5}]$ in $4841$ steps. The barrier of Polymath15 is near $X=6\times10^{10}$, height $3\times10^{10}$ for $\xi_{0.2}$, and reaching it in this framework requires the evaluator of the programme below.
